\chapter{Ongoing and future work}
\label{ch:future}

\section{Running and prepared experiments}
\label{sec:future-running}

\begin{itemize}
  \item \textbf{A fair 7B run.} Qwen2.5-7B with a lower learning rate
        ($5 \times 10^{-5}$), so that model size is no longer confounded with unstable training.
  \item \textbf{Meaning without numbers.} The \emph{rich} format did not help
        (Section~\ref{sec:results-format}). The \emph{semantic only} format, which gives the band
        and the phrase without any number, is implemented and will show how much the model can
        learn from the meaning alone.
\end{itemize}

\section{Clinically motivated features}

The TI-RADS-motivated features are the most promising direction (Section~\ref{sec:results-clinical}).
Three steps follow:

\begin{itemize}
  \item \textbf{Echogenicity of the solid part.} The echogenicity ratio is strongly tied to the
        anechoic fraction (correlation $-0.78$): it measures fluid rather than the echogenicity of
        tissue. Computed on the non-fluid pixels only, as TI-RADS prescribes, it may point the
        expected way.
  \item \textbf{Their meaning for the language model.} The plain-language formats did not help
        with radiomics features, whose names mean little to a language model. Phrases such as
        ``taller than wide'' or ``a more lobulated outline'' are the language of radiology reports;
        if meaning helps anywhere, it should help here. A prompt with the TI-RADS features alone
        isolates this question.
  \item \textbf{A pre-specified ordering for Learn-then-Test.} Ordering the candidate thresholds
        on validation certified an auto-malignant rule for run I (Section~\ref{sec:results-ltt}).
        Fixed in advance, the rule can be tested on new seeds and splits.
\end{itemize}

\section{Robustness}

The best setting will be repeated with two or three seeds and reported as a mean with a range,
so that differences between settings can be judged against the variation between runs.

\section{A second dataset}

ThyroidXL will be used as a separate dataset, not merged with TN3K. It can show whether the
results hold on other data and, with more calibration nodules, whether the triage guarantee
becomes reachable.
